\documentclass[5p,twocolumn,times]{elsarticle}

\usepackage[utf8]{inputenc}
\usepackage[T1]{fontenc}
\usepackage[english]{babel}
\usepackage{hyperref}

\journal{HealthAI 2026}

\hypersetup{
  colorlinks=true,
  linkcolor=black,
  citecolor=black,
  urlcolor=blue
}

\begin{document}

\begin{frontmatter}

\title{DentIA: An Artificial Intelligence-Based Decision Support Tool for Dental Implant Planning Through Biomechanical Simulation and Knowledge-Grounded Clinical Assistance}

\author[ust]{Cesar Hernando Valencia Niño}
\ead{cesar.valencia@ustabuca.edu.co}
\ead[url]{https://orcid.org/0000-0001-6077-6458}
\address[ust]{Facultad de Ingeniería Mecatrónica, Universidad Santo Tomás - Seccional Bucaramanga, Colombia}

\begin{abstract}
Dental implant planning requires the integration of anatomical, prosthetic, biomechanical, and patient-related factors. In full-arch oral rehabilitation, decisions are shaped by bone quality, implant distribution and angulation, cantilever length, bruxism, occlusal load, bar material, and prosthetic design. Because these variables interact with one another, they can be difficult to analyze simultaneously during clinical case discussion or during undergraduate and postgraduate training. This work presents DentIA, a web-based decision support tool developed to assist implant dentistry planning by combining a biomechanical simulator with an artificial intelligence agent connected to a structured clinical knowledge base.

The biomechanical simulator allows users to explore implant-prosthetic scenarios by adjusting clinically relevant variables, including arch type, bone density, implant length and diameter, implant configuration, cantilever extension, bar profile, and material selection. From these inputs, the platform provides comparative guidance on relative biomechanical risk and helps users visualize how each factor may influence prosthetic planning. The simulator is intended to support analytical reasoning and case discussion, rather than to generate definitive treatment prescriptions.

The second component is an AI-powered clinical assistant that enables natural-language consultation of concepts related to implant selection, support bar design, material properties, rehabilitation protocols, biomechanical principles, and risk factors. By linking the agent to a curated knowledge base, DentIA aims to provide context-aware explanations and support evidence-informed discussion of clinical cases.

DentIA was conceived as an educational and clinical support platform, not as a substitute for professional diagnosis, specialist judgment, or individualized treatment planning. Its main contribution lies in bringing together simulation-based analysis and knowledge-grounded AI assistance within a single digital workflow. This integrated approach may help students, educators, and clinicians organize information, compare alternatives, and strengthen the discussion of implant rehabilitation cases. Future work will include expert validation, usability assessment, and comparison between system outputs and specialist recommendations.
\end{abstract}

\begin{keyword}
artificial intelligence \sep dental implants \sep clinical decision support \sep biomechanical simulation \sep implant dentistry \sep digital health
\end{keyword}

\end{frontmatter}

\section{Introduction}

% Add the clinical context, problem statement, and rationale for developing DentIA.

\section{Materials and Methods}

% Describe the design of the web platform, software architecture, knowledge base,
% biomechanical simulator, AI agent, variables, and development workflow.

\section{Results}

% Present the implemented modules, interface, example clinical scenario,
% simulator outputs, and AI-assisted consultation workflow.

\section{Discussion}

% Discuss clinical and educational value, comparison with existing digital tools,
% limitations, safety boundaries, and future validation.

\section{Conclusion}

% Summarize the contribution of DentIA and the next steps for validation.

\section*{Acknowledgments}

% Add acknowledgments if applicable.

\section*{Funding}

% Add funding information if applicable.

\section*{Conflicts of Interest}

% Add conflict of interest statement.

\section*{Data Availability}

% Add data availability statement if applicable.

\section*{References}

% Add references here using the journal or conference citation style.

\end{document}
